% !TeX TXS-program:bibliography = txs:///biber
\documentclass[14pt, russian]{scrartcl}
\let\counterwithout\relax
\let\counterwithin\relax
\usepackage{float}
\usepackage{xcolor}
\usepackage{extsizes}
\usepackage{subfig}
\usepackage[export]{adjustbox}
\usepackage{tocvsec2} % возможность менять учитываемую глубину разделов в оглавлении
\usepackage[subfigure]{tocloft}
% \usepackage[newfloat]{minted}  % Отключено для ускорения компиляции в Overleaf
% \captionsetup[listing]{position=top}

% \makeatletter
% \AddToHook{begindocument/before}{\@ifpackageloaded{minted}{\removefromtoclist[float]{lol}}{}}
% \makeatother

\AtBeginEnvironment{figure}{\vspace{0.5cm}}
\AtBeginEnvironment{table}{\vspace{0.5cm}}
\AtBeginEnvironment{listing}{\vspace{0.5cm}}
\AtBeginEnvironment{algorithm}{\vspace{0.5cm}}
% \AtBeginEnvironment{minted}{\vspace{-0.5cm}}  % Отключено вместе с minted

\usepackage{fancyvrb}
\usepackage[normalem]{ulem}
\usepackage{bm,mathrsfs,ifsym} %зачеркивания, особо жирный стиль и RSFS начертание
\usepackage{sectsty} % переопределение стилей подразделов
%%%%%%%%%%%%%%%%%%%%%%%

%%% Поля и разметка страницы %%%
\usepackage{pdflscape}                              % Для включения альбомных страниц
\usepackage{geometry}                               % Для последующего задания полей
\geometry{a4paper,tmargin=2cm,bmargin=2cm,lmargin=3cm,rmargin=1cm} % тоже самое, но лучше

%%% Математические пакеты %%%
\usepackage{amsthm,amsfonts,amsmath,amssymb,amscd}  % Математические дополнения от AMS
\usepackage{mathtools}                              % Добавляет окружение multlined
\usepackage[perpage]{footmisc}
%\usepackage{times}

%%%% Установки для размера шрифта 14 pt %%%%
%% Формирование переменных и констант для сравнения (один раз для всех подключаемых файлов)%%
%% должно располагаться до вызова пакета fontspec или polyglossia, потому что они сбивают его работу
%\newlength{\curtextsize}
%\newlength{\bigtextsize}
%\setlength{\bigtextsize}{13pt}
\KOMAoptions{fontsize=14pt}

\makeatletter
\def\showfontsize{\f@size{} point}
\makeatother

\makeatletter
\setlength{\@fptop}{0pt}
\makeatother

%\makeatletter
%\show\f@size                                       % неплохо для отслеживания, но вызывает стопорение процесса, если документ компилируется без команды  -interaction=nonstopmode 
%\setlength{\curtextsize}{\f@size pt}
%\makeatother

%шрифт Times New Roman
\usepackage{times}  % Пакет для шрифта Times New Roman

   %%% Решение проблемы копирования текста в буфер кракозябрами
%    \input glyphtounicode.tex
%    \input glyphtounicode-cmr.tex %from pdfx package
%    \pdfgentounicode=1
    \usepackage{cmap}                               % Улучшенный поиск русских слов в полученном pdf-файле
    \usepackage[T1]{fontenc}                       % Поддержка русских букв
    \usepackage[utf8]{inputenc}                     % Кодировка utf8
    \usepackage[english, main=russian]{babel}            % Языки: русский, английский
%   \IfFileExists{pscyr.sty}{\usepackage{pscyr}}{}  % Красивые русские шрифты
%\renewcommand{\rmdefault}{ftm}
%%% Оформление абзацев %%%
\usepackage{indentfirst}                            % Красная строка
%\usepackage{eskdpz}

%%% Таблицы %%%
\usepackage{longtable}                              % Длинные таблицы
\usepackage{multirow,makecell,array}                % Улучшенное форматирование таблиц
\usepackage{booktabs}                               % Возможность оформления таблиц в классическом книжном стиле (при правильном использовании не противоречит ГОСТ)

%%% Общее форматирование
\usepackage{soulutf8}                               % Поддержка переносоустойчивых подчёркиваний и зачёркиваний
\usepackage{icomma}                                 % Запятая в десятичных дробях



%%% Изображения %%%
\usepackage{graphicx}                               % Подключаем пакет работы с графикой
\usepackage{wrapfig}

%%% Списки %%%
\usepackage{enumitem}

%%% Подписи %%%
\usepackage{caption}                                % Для управления подписями (рисунков и таблиц) % Может управлять номерами рисунков и таблиц с caption %Иногда может управлять заголовками в списках рисунков и таблиц
%% Использование:
%\begin{table}[h!]\ContinuedFloat - чтобы не переключать счетчик
%\captionsetup{labelformat=continued}% должен стоять до самого caption
%\caption{}
% либо ручками \caption*{Продолжение таблицы~\ref{...}.} :)

%%% Интервалы %%%
\addto\captionsrussian{%
  \renewcommand{\listingname}{Листинг}%
}
%%% Счётчики %%%
\usepackage[figure,table,section]{totalcount}               % Счётчик рисунков и таблиц
\DeclareTotalCounter{lstlisting}
\usepackage{totcount}                               % Пакет создания счётчиков на основе последнего номера подсчитываемого элемента (может требовать дважды компилировать документ)
\usepackage{totpages}                               % Счётчик страниц, совместимый с hyperref (ссылается на номер последней страницы). Желательно ставить последним пакетом в преамбуле

%%% Продвинутое управление групповыми ссылками (пока только формулами) %%%
%% Кодировки и шрифты %%%

%   \newfontfamily{\cyrillicfont}{Times New Roman}
%   \newfontfamily{\cyrillicfonttt}{CMU Typewriter Text}
	%\setmainfont{Times New Roman}
	%\newfontfamily\cyrillicfont{Times New Roman}
	%\setsansfont{Times New Roman}                    %% задаёт шрифт без засечек
%	\setmonofont{Liberation Mono}               %% задаёт моноширинный шрифт
%    \IfFileExists{pscyr.sty}{\renewcommand{\rmdefault}{ftm}}{}
%%% Интервалы %%%
%linespread-реализация ближе к реализации полуторного интервала в ворде.
%setspace реализация заточена под шрифты 10, 11, 12pt, под остальные кегли хуже, но всё же ближе к типографской классике. 
\linespread{1.3}                    % Полуторный интервал (ГОСТ Р 7.0.11-2011, 5.3.6)
%\renewcommand{\@biblabel}[1]{#1}

%%% Гиперссылки %%%
\usepackage{hyperref}

%%% Выравнивание и переносы %%%
\sloppy                             % Избавляемся от переполнений
\clubpenalty=10000                  % Запрещаем разрыв страницы после первой строки абзаца
\widowpenalty=10000                 % Запрещаем разрыв страницы после последней строки абзаца

\makeatletter % малые заглавные, small caps shape
\let\@@scshape=\scshape
\renewcommand{\scshape}{%
  \ifnum\strcmp{\f@series}{bx}=\z@
    \usefont{T1}{cmr}{bx}{sc}%
  \else
    \ifnum\strcmp{\f@shape}{it}=\z@
      \fontshape{scsl}\selectfont
    \else
      \@@scshape
    \fi
  \fi}
\makeatother

%%% Подписи %%%
%\captionsetup{%
%singlelinecheck=off,                % Многострочные подписи, например у таблиц
%skip=2pt,                           % Вертикальная отбивка между подписью и содержимым рисунка или таблицы определяется ключом
%justification=centering,            % Центрирование подписей, заданных командой \caption
%}
%%%        Подключение пакетов                 %%%
\usepackage{ifthen}                 % добавляет ifthenelse
%%% Инициализирование переменных, не трогать!  %%%
\newcounter{intvl}
\newcounter{otstup}
\newcounter{contnumeq}
\newcounter{contnumfig}
\newcounter{contnumtab}
\newcounter{pgnum}
\newcounter{bibliosel}
\newcounter{chapstyle}
\newcounter{headingdelim}
\newcounter{headingalign}
\newcounter{headingsize}
\newcounter{tabcap}
\newcounter{tablaba}
\newcounter{tabtita}
%%%%%%%%%%%%%%%%%%%%%%%%%%%%%%%%%%%%%%%%%%%%%%%%%%

%%% Область упрощённого управления оформлением %%%

%% Интервал между заголовками и между заголовком и текстом
% Заголовки отделяют от текста сверху и снизу тремя интервалами (ГОСТ Р 7.0.11-2011, 5.3.5)
\setcounter{intvl}{3}               % Коэффициент кратности к размеру шрифта

%% Отступы у заголовков в тексте
\setcounter{otstup}{0}              % 0 --- без отступа; 1 --- абзацный отступ

%% Нумерация формул, таблиц и рисунков
\setcounter{contnumeq}{1}           % Нумерация формул: 0 --- пораздельно (во введении подряд, без номера раздела); 1 --- сквозная нумерация по всей диссертации
\setcounter{contnumfig}{1}          % Нумерация рисунков: 0 --- пораздельно (во введении подряд, без номера раздела); 1 --- сквозная нумерация по всей диссертации
\setcounter{contnumtab}{1}          % Нумерация таблиц: 0 --- пораздельно (во введении подряд, без номера раздела); 1 --- сквозная нумерация по всей диссертации

%% Оглавление
\setcounter{pgnum}{0}               % 0 --- номера страниц никак не обозначены; 1 --- Стр. над номерами страниц (дважды компилировать после изменения)

%% Библиография
\setcounter{bibliosel}{1}           % 0 --- встроенная реализация с загрузкой файла через движок bibtex8; 1 --- реализация пакетом biblatex через движок biber

%% Текст и форматирование заголовков
\setcounter{chapstyle}{1}           % 0 --- разделы только под номером; 1 --- разделы с названием "Глава" перед номером
\setcounter{headingdelim}{1}        % 0 --- номер отделен пропуском в 1em или \quad; 1 --- номера разделов и приложений отделены точкой с пробелом, подразделы пропуском без точки; 2 --- номера разделов, подразделов и приложений отделены точкой с пробелом.

%% Выравнивание заголовков в тексте
\setcounter{headingalign}{0}        % 0 --- по центру; 1 --- по левому краю

%% Размеры заголовков в тексте
\setcounter{headingsize}{0}         % 0 --- по ГОСТ, все всегда 14 пт; 1 --- пропорционально изменяющийся размер в зависимости от базового шрифта

%% Подпись таблиц
\setcounter{tabcap}{0}              % 0 --- по ГОСТ, номер таблицы и название разделены тире, выровнены по левому краю, при необходимости на нескольких строках; 1 --- подпись таблицы не по ГОСТ, на двух и более строках, дальнейшие настройки: 
%Выравнивание первой строки, с подписью и номером
\setcounter{tablaba}{2}             % 0 --- по левому краю; 1 --- по центру; 2 --- по правому краю
%Выравнивание строк с самим названием таблицы
\setcounter{tabtita}{1}             % 0 --- по левому краю; 1 --- по центру; 2 --- по правому краю

%%% Рисунки %%%
\DeclareCaptionLabelSeparator*{emdash}{~--- }             % (ГОСТ 2.105, 4.3.1)
\captionsetup[figure]{labelsep=emdash,font=onehalfspacing,position=bottom}

%%% Таблицы %%%
\ifthenelse{\equal{\thetabcap}{0}}{%
    \newcommand{\tabcapalign}{\raggedright}  % по левому краю страницы или аналога parbox
}

\ifthenelse{\equal{\thetablaba}{0} \AND \equal{\thetabcap}{1}}{%
    \newcommand{\tabcapalign}{\raggedright}  % по левому краю страницы или аналога parbox
}

\ifthenelse{\equal{\thetablaba}{1} \AND \equal{\thetabcap}{1}}{%
    \newcommand{\tabcapalign}{\centering}    % по центру страницы или аналога parbox
}

\ifthenelse{\equal{\thetablaba}{2} \AND \equal{\thetabcap}{1}}{%
    \newcommand{\tabcapalign}{\raggedleft}   % по правому краю страницы или аналога parbox
}

\ifthenelse{\equal{\thetabtita}{0} \AND \equal{\thetabcap}{1}}{%
    \newcommand{\tabtitalign}{\raggedright}  % по левому краю страницы или аналога parbox
}

\ifthenelse{\equal{\thetabtita}{1} \AND \equal{\thetabcap}{1}}{%
    \newcommand{\tabtitalign}{\centering}    % по центру страницы или аналога parbox
}

\ifthenelse{\equal{\thetabtita}{2} \AND \equal{\thetabcap}{1}}{%
    \newcommand{\tabtitalign}{\raggedleft}   % по правому краю страницы или аналога parbox
}

\DeclareCaptionFormat{tablenocaption}{\tabcapalign #1\strut}        % Наименование таблицы отсутствует
\ifthenelse{\equal{\thetabcap}{0}}{%
    \DeclareCaptionFormat{tablecaption}{\tabcapalign #1#2#3}
    \captionsetup[table]{labelsep=emdash}                       % тире как разделитель идентификатора с номером от наименования
}{%
    \DeclareCaptionFormat{tablecaption}{\tabcapalign #1#2\par%  % Идентификатор таблицы на отдельной строке
        \tabtitalign{#3}}                                       % Наименование таблицы строкой ниже
    \captionsetup[table]{labelsep=space}                        % пробельный разделитель идентификатора с номером от наименования
}
\captionsetup[table]{format=tablecaption,singlelinecheck=off,font=onehalfspacing,position=top,skip=-5pt}  % многострочные наименования и прочее
\DeclareCaptionLabelFormat{continued}{Продолжение таблицы~#2}
\setlength{\belowcaptionskip}{.2cm}
\setlength{\intextsep}{0ex}

%%% Подписи подрисунков %%%
\renewcommand{\thesubfigure}{\asbuk{subfigure}}           % Буквенные номера подрисунков
\captionsetup[subfigure]{font={normalsize},               % Шрифт подписи названий подрисунков (не отличается от основного)
    labelformat=brace,                                    % Формат обозначения подрисунка
    justification=centering,                              % Выключка подписей (форматирование), один из вариантов            
}
%\DeclareCaptionFont{font12pt}{\fontsize{12pt}{13pt}\selectfont} % объявляем шрифт 12pt для использования в подписях, тут же надо интерлиньяж объявлять, если не наследуется
%\captionsetup[subfigure]{font={font12pt}}                 % Шрифт подписи названий подрисунков (всегда 12pt)

%%% Настройки гиперссылок %%%

\definecolor{linkcolor}{rgb}{0.0,0,0}
\definecolor{citecolor}{rgb}{0,0.0,0}
\definecolor{urlcolor}{rgb}{0,0,0}

\hypersetup{
    linktocpage=true,           % ссылки с номера страницы в оглавлении, списке таблиц и списке рисунков
%    linktoc=all,                % both the section and page part are links
%    pdfpagelabels=false,        % set PDF page labels (true|false)
    plainpages=true,           % Forces page anchors to be named by the Arabic form  of the page number, rather than the formatted form
    colorlinks,                 % ссылки отображаются раскрашенным текстом, а не раскрашенным прямоугольником, вокруг текста
    linkcolor={linkcolor},      % цвет ссылок типа ref, eqref и подобных
    citecolor={citecolor},      % цвет ссылок-цитат
    urlcolor={urlcolor},        % цвет гиперссылок
    pdflang={ru},
}
\urlstyle{same}
%%% Шаблон %%%
%\DeclareRobustCommand{\todo}{\textcolor{red}}       % решаем проблему превращения названия цвета в результате \MakeUppercase, http://tex.stackexchange.com/a/187930/79756 , \DeclareRobustCommand protects \todo from expanding inside \MakeUppercase
\setlength{\parindent}{2.5em}                       % Абзацный отступ. Должен быть одинаковым по всему тексту и равен пяти знакам (ГОСТ Р 7.0.11-2011, 5.3.7).

%%% Списки %%%
% Используем дефис для ненумерованных списков (ГОСТ 2.105-95, 4.1.7)
%\renewcommand{\labelitemi}{\normalfont\bfseries~{---}} 
\renewcommand{\labelitemi}{\bfseries~{---}} 
\setlist{nosep,%                                    % Единый стиль для всех списков (пакет enumitem), без дополнительных интервалов.
    labelindent=\parindent,leftmargin=*%            % Каждый пункт, подпункт и перечисление записывают с абзацного отступа (ГОСТ 2.105-95, 4.1.8)
}
%%%%%%%%%%%%%%%%%%%%%%
%\usepackage{xltxtra} % load xunicode

\usepackage{ragged2e}
\usepackage[explicit]{titlesec}
\usepackage{placeins}
\usepackage{xparse}
\usepackage{csquotes}

\usepackage{listingsutf8}
\usepackage{url} %пакеты расширений
\usepackage{algorithm, algorithmicx}
\usepackage[noend]{algpseudocode}
\usepackage{blkarray}
\usepackage{chngcntr}
\usepackage{tabularx}
\usepackage[backend=biber, 
    bibstyle=gost-numeric,
    citestyle=nature]{biblatex}
\newcommand*\template[1]{\text{<}#1\text{>}}
  
\titleformat{name=\section,numberless}[block]{\normalfont\large\centering}{}{0em}{#1}
\titleformat{\section}[block]{\normalfont\large\bfseries\raggedright}{}{0em}{\thesection\hspace{0.25em}#1}
\titleformat{\subsection}[block]{\normalfont\normalsize\bfseries\raggedright}{}{0em}{\thesubsection\hspace{0.25em}#1}
\titleformat{\subsubsection}[block]{\normalfont\normalsize\bfseries\raggedright}{}{0em}{\thesubsubsection\hspace{0.25em}#1}

\let\Algorithm\algorithm
\renewcommand\algorithm[1][]{\Algorithm[#1]\setstretch{1.5}}
%\renewcommand{\listingscaption}{Листинг}

\usepackage{pifont}
\usepackage{calc}
\usepackage{suffix}
\usepackage{csquotes}
\DeclareQuoteStyle{russian}
    {\guillemotleft}{\guillemotright}[0.025em]
    {\quotedblbase}{\textquotedblleft}
\ExecuteQuoteOptions{style=russian}
\newcommand{\enq}[1]{\enquote{#1}}  
\newcommand{\eng}[1]{\begin{english}#1\end{english}}
% Подчиненные счетчики в окружениях http://old.kpfu.ru/journals/izv_vuz/arch/sample1251.tex
\newcounter{cTheorem} 
\newcounter{cDefinition}
\newcounter{cConsequent}
\newcounter{cExample}
\newcounter{cLemma}
\newcounter{cConjecture}
\newtheorem{Theorem}{Теорема}[cTheorem]
\newtheorem{Definition}{Определение}[cDefinition]
\newtheorem{Consequent}{Следствие}[cConsequent]
\newtheorem{Example}{Пример}[cExample]
\newtheorem{Lemma}{Лемма}[cLemma]
\newtheorem{Conjecture}{Гипотеза}[cConjecture]

\renewcommand{\theTheorem}{\arabic{Theorem}}
\renewcommand{\theDefinition}{\arabic{Definition}}
\renewcommand{\theConsequent}{\arabic{Consequent}}
\renewcommand{\theExample}{\arabic{Example}}
\renewcommand{\theLemma}{\arabic{Lemma}}
\renewcommand{\theConjecture}{\arabic{Conjecture}}
%\makeatletter
\NewDocumentCommand{\Newline}{}{\text{\\}}
\newcommand{\sequence}[2]{\ensuremath \left(#1,\ \dots,\ #2\right)}

\definecolor{mygreen}{rgb}{0,0.6,0}
\definecolor{mygray}{rgb}{0.5,0.5,0.5}
\definecolor{mymauve}{rgb}{0.58,0,0.82}
\renewcommand{\listalgorithmname}{Список алгоритмов}
\floatname{algorithm}{Листинг}
\renewcommand{\lstlistingname}{Листинг}
\renewcommand{\thealgorithm}{\arabic{algorithm}}

\newcommand{\refAlgo}[1]{(листинг \ref{#1})}
\newcommand{\refImage}[1]{(рисунок \ref{#1})}

\renewcommand{\theenumi}{\arabic{enumi}.}% Меняем везде перечисления на цифра.цифра	
\renewcommand{\labelenumi}{\arabic{enumi}.}% Меняем везде перечисления на цифра.цифра
\renewcommand{\theenumii}{\arabic{enumii}}% Меняем везде перечисления на цифра.цифра
\renewcommand{\labelenumii}{(\arabic{enumii})}% Меняем везде перечисления на цифра.цифра
\renewcommand{\theenumiii}{\roman{enumiii}}% Меняем везде перечисления на цифра.цифра
\renewcommand{\labelenumiii}{(\roman{enumiii})}% Меняем везде перечисления на цифра.цифра
\renewcommand{\labelitemi}{---}
\renewcommand{\labelitemii}{---}

\makeatletter
\def\p@subsection{}
\def\p@subsubsection{\thesection\,\thesubsection\,}
\makeatother
\newcommand{\anonsection}[1]{\cleardoublepage
\phantomsection
\addcontentsline{toc}{section}{\protect\numberline{}#1}
\section*{#1}\vspace*{2.5ex} % По госту положены 3 пустые строки после заголовка ненумеруемого раздела
}
\newcommand{\sectionbreak}{\clearpage}
\renewcommand{\sectionfont}{\normalsize} % Сбиваем стиль оглавления в стандартный
\renewcommand{\cftsecleader}{\cftdotfill{\cftdotsep}} % Точки в оглавлении напротив разделов

\renewcommand{\cftsecfont}{\normalfont\large} % Переключение на times в содержании
\renewcommand{\cftsubsecfont}{\normalfont\large} % Переключение на times в содержании

\setlength{\cftsecindent}{0pt}% Убираем отступы в содержании для \section
\setlength{\cftsubsecindent}{0pt}% Убираем отступы в содержании для \subsection
\setlength{\cftsubsubsecindent}{0pt}% Убираем отступы в содержании для \subsubsection

\usepackage{caption} 
%\captionsetup[table]{justification=raggedleft} 
%\captionsetup[figure]{justification=centering,labelsep=endash}
\usepackage{amsmath}    % \bar    (матрицы и проч. ...)
\usepackage{amsfonts}   % \mathbb (символ для множества действительных чисел и проч. ...)
\usepackage{mathtools}  % \abs, \norm
    \DeclarePairedDelimiter\abs{\lvert}{\rvert} % операция модуля
    \DeclarePairedDelimiter\norm{\lVert}{\rVert} % операция нормы
\DeclareTextCommandDefault{\textvisiblespace}{%
  \mbox{\kern.06em\vrule \@height.3ex}%
  \vbox{\hrule \@width.3em}%
  \hbox{\vrule \@height.3ex}}    
\newsavebox{\spacebox}
\begin{lrbox}{\spacebox}
\verb*! !
\end{lrbox}
\newcommand{\aspace}{\usebox{\spacebox}}
\DeclareTotalCounter{listing}
\def\figurename{Рисунок}

\makeatletter
\renewcommand*{\p@subsubsection}{}
\makeatother

\newcommand{\bmstuheader}
{\vspace*{-30pt}
\hspace{-45pt}
\begin{minipage}{0.17\textwidth}
\hspace*{-20pt}\centering
\includegraphics[width=1.3\textwidth]{emblem.png}
\end{minipage}
\begin{minipage}{0.82\textwidth}\small \textbf{
\vspace*{-0.7ex}
\hspace*{-10pt}\centerline{Министерство науки и высшего образования Российской Федерации}
\vspace*{-0.7ex}
\centerline{Федеральное государственное автономное образовательное учреждение }
\vspace*{-0.7ex}
\centerline{высшего образования}
\vspace*{-0.7ex}
\centerline{<<Московский государственный технический университет}
\vspace*{-0.7ex}
\centerline{имени Н.Э. Баумана}
\vspace*{-0.7ex}
\centerline{(национальный исследовательский университет)>>}
\vspace*{-0.7ex}
\centerline{(МГТУ им. Н.Э. Баумана)}}
\end{minipage}

\vspace{-2pt}
\hspace{-34.5pt}\rule{\textwidth}{0.5pt}

\vspace*{-18.3pt}
\hspace{-34.5pt}\rule{\textwidth}{2.5pt}
 
\vspace{0.5ex}
\noindent \small ФАКУЛЬТЕТ\hspace{80pt} <<Информатика и системы управления>>

\vspace*{-16pt}
\hspace{35pt}\rule{0.855\textwidth}{0.4pt}

\vspace{0.5ex}
\noindent \small КАФЕДРА\hspace{50pt} <<Теоретическая информатика и компьютерные технологии>>

\vspace*{-16pt}
\hspace{25pt}\rule{0.875\textwidth}{0.4pt}
}
\addbibresource{biblio.bib}

\begin{document}
\sloppy

\begin{titlepage}
\thispagestyle{empty}

\bmstuheader
 
\vspace{3em}
 
\begin{center}
\Large \textbf{РАСЧЕТНО-ПОЯСНИТЕЛЬНАЯ ЗАПИСКА\\\textbf{\textit{К КУРСОВОЙ РАБОТЕ\\НА ТЕМУ:}} \\}
\end{center}

\vspace*{-6ex} 
\begin{center}
\Large{\textit{\textbf{<<Компилятор подмножества Фортрана>>}}}

\vspace*{-3ex}
\rule{0.9\textwidth}{1.2pt}

\vspace*{-0.2ex}
\rule{0.9\textwidth}{1.2pt}

\vspace*{-0.2ex}
\rule{0.9\textwidth}{1.2pt}

\vspace*{-0.2ex}
\rule{0.9\textwidth}{1.2pt}

\vspace*{-0.2ex}
\rule{0.9\textwidth}{1.2pt}
\end{center}
 
\vspace{\fill}

\newlength{\ML}
\settowidth{\ML}{«\underline{\hspace{0.7cm}}» \underline{\hspace{2cm}}}

\noindent Студент \underline{\hspace{1.5cm}} \hfill \underline{\hspace{4cm}}\quad
\underline{\hspace{4cm}}

\vspace{-2.1ex}
\noindent\hspace{9ex}\scriptsize{(Группа)}\normalsize\hspace{170pt}\hspace{2ex}\scriptsize{(Подпись, дата)}\normalsize\hspace{30pt}\hspace{6ex}\scriptsize{(И.О. Фамилия)}\normalsize

\bigskip

\noindent Руководитель  \hfill \underline{\hspace{4cm}}\quad
\underline{\hspace{4cm}}

\vspace{-2ex}
\noindent\hspace{13.5ex}\normalsize\hspace{170pt}\hspace{2ex}\scriptsize{(Подпись, дата)}\normalsize\hspace{30pt}\hspace{6ex}\scriptsize{(И.О. Фамилия)}\normalsize
\vfill

\begin{center}
\textsl{\the\year{} г.}
\end{center}
\end{titlepage}

\setlength{\tabcolsep}{3pt}
\newpage
\setcounter{page}{2}

\renewcommand\contentsname{\hfill{\normalfont{СОДЕРЖАНИЕ}}\hfill}
\setcounter{tocdepth}{3}
\tableofcontents
\newpage
\anonsection{ВВЕДЕНИЕ}

Эффективность программных комплексов для задач вычислительной математики, физического моделирования и обработки больших данных во многом определяется удачной организацией исполнения алгоритмов на современной вычислительной архитектуре. Для языка Фортран, который традиционно занимает ведущие позиции в сфере высокопроизводительных научных вычислений, эта зависимость проявляется особенно ярко. Значительная часть таких прикладных программ строится вокруг интенсивной обработки многомерных массивов данных, заключенных в многократно вложенные циклические конструкции. Именно эти фрагменты кода задают основную вычислительную нагрузку, формируют профиль обращений к иерархической памяти и служат главным ограничивающим фактором для распараллеливания \cite{aho2007dragon}.

За синтаксической простотой циклов скрывается сложнейший механизм обмена данными между процессорными регистрами, многоуровневой кэш-памятью и оперативной памятью. В условиях так называемой проблемы «стены памяти» (memory wall), когда скорость арифметических операций многократно превосходит пропускную способность шины памяти, существенный выигрыш в производительности достигается только в том случае, когда транслятор глубоко анализирует зависимости между итерациями и радикально изменяет порядок обхода пространства итераций. Это позволяет повысить пространственную и временную локальность данных, а также выделить независимые блоки вычислений для параллельного исполнения \cite{muchnick1997}.

Для подмножества языка Фортран автоматическое применение подобных полиэдральных трансформаций концептуально разрешимо и математически обосновано. Строгая семантика массивов и отсутствие сложного наложения адресов (алиасинга) облегчают статический анализ, позволяя оптимизирующему компилятору уверенно идентифицировать регулярные вычислительные ядра и доказывать формальную корректность вносимых изменений.

Особый научный и практический интерес представляют гнёзда циклов итерационного типа, типичные для конечно-разностных схем, методов релаксации и клеточных автоматов. В таких конструкциях внешний цикл отвечает за смену временных слоев или шагов сходимости, а внутренние реализуют обход элементов пространственной сетки по заданному шаблону. Для них возможно построение формальной алгебраической модели пространства итераций и применение упорядоченной цепочки преобразований: скашивания, блочного разбиения (тайлинга) и метода гиперплоскостей \cite{metelica2024}.

\textbf{Цель данной работы} состоит в детальном теоретическом обосновании методов оптимизации многомерных вложенных циклов и их адаптации для применения в архитектуре оптимизирующего компилятора подмножества языка Фортран.

Для достижения поставленной цели были сформулированы следующие \textbf{задачи}:
\begin{itemize}
    \item Проанализировать архитектурные предпосылки и теоретические основы цикловых оптимизаций, включая формальные методы описания пространств итераций.
    \item Изучить базовые полиэдральные преобразования циклов: перестановку, скашивание, блочное разбиение и метод гиперплоскостей.
    \item Разработать и обосновать комплексный алгоритм ускорения гнёзд циклов итерационного типа.
    \item Сформулировать аналитическую модель для определения оптимальных размеров блоков с учетом аппаратных характеристик кэш-памяти.
\end{itemize}

\textbf{Объектом исследования} выступают методы статического анализа и эквивалентных алгебраических преобразований исходного кода высокопроизводительных численных программ.

\textbf{Предметом исследования} являются автоматические преобразования гнёзд циклов итерационного типа, направленные на минимизацию промахов кэш-памяти и извлечение параллелизма на уровне потоков.

\section{Теоретические основы оптимизации циклических вычислений}

\subsection{Архитектурные предпосылки оптимизации и проблема памяти}

Историческое развитие вычислительной техники характеризуется неравномерным ростом производительности различных компонентов системы. В то время как вычислительная мощность процессорных ядер (выражаемая в количестве выполняемых инструкций за такт) росла экспоненциально согласно закону Мура, пропускная способность и, что более важно, латентность доступа к основной оперативной памяти улучшались значительно более скромными темпами. Этот нарастающий разрыв привел к феномену, известному в компьютерной инженерии как «стена памяти» (memory wall). 

В современных процессорных архитектурах выполнение простейшей арифметической операции, такой как сложение двух чисел с плавающей точкой, может занимать всего один такт процессорного времени, тогда как извлечение операндов для этой операции из основной оперативной памяти требует сотен тактов простоя (stall cycles). Для смягчения этого дисбаланса современные процессоры оснащаются многоуровневой иерархией кэш-памяти (обычно L1, L2 и L3), где каждый последующий уровень обладает большей емкостью, но меньшей скоростью доступа \cite{allenKennedy2002}.

Однако наличие кэш-памяти само по себе не гарантирует высокой производительности. Эффективность кэширования всецело зависит от свойств локальности данных в исполняемой программе. Если программа обращается к памяти хаотично или с большими шагами, кэш-память постоянно обновляется (trashing), и процессор вынужден простаивать в ожидании данных. Таким образом, одной из главнейших задач современного оптимизирующего компилятора становится трансформация исходного кода программы таким образом, чтобы максимизировать использование быстрых уровней памяти. В программах для научных вычислений подавляющая часть времени затрачивается на обработку массивов именно внутри циклов. Они задают порядок обхода данных, что напрямую влияет на пространственную и временную локальность \cite{cooperTorczon2011}.

\subsection{Циклы как основной объект анализа компилятора}

С точки зрения компилятора, цикл представляет собой конструкцию управления, задающую многократное повторение блока инструкций. Для программ на языке Фортран наиболее характерны счетные циклы (в Фортране -- оператор `DO`), количество итераций которых может быть определено до начала выполнения самого тела цикла. Это свойство кардинально упрощает статический анализ потока управления (Control Flow Analysis) и потока данных (Data Flow Analysis).

Наиболее удобными и благодарными объектами для агрессивной оптимизации являются так называемые тесные гнёзда циклов (perfectly nested loops). В тесном гнезде все операторы присваивания и вычисления сосредоточены исключительно в самом внутреннем цикле, а между заголовками циклов отсутствуют какие-либо инструкции. Такая строгая структура позволяет рассматривать всё гнездо как единый многомерный объект и применять к нему сложный математический аппарат полиэдральной модели \cite{quillere2000}. Строгая семантика массивов в Фортране, запрещающая неявное наложение адресов (алиасинг) в параметрах подпрограмм по умолчанию, позволяет компилятору уверенно делать выводы о независимости ячеек памяти и безопасно применять масштабные трансформации \cite{feautrier1992affine}.

\subsection{Информационные зависимости по данным}

Краеугольным камнем теории распараллеливания и цикловых преобразований является концепция информационных зависимостей (data dependencies). Критерием корректности абсолютно любого преобразования компилятора является сохранение семантики исходной программы, что на практике сводится к сохранению порядка выполнения зависимых инструкций \cite{banerjee1988}.

Информационная зависимость возникает между двумя обращениями к памяти, если они потенциально могут обращаться к одной и той же ячейке, и хотя бы одно из этих обращений является операцией записи (модификацией). Формально выделяют три фундаментальных типа зависимостей:
\begin{enumerate}
    \item \textbf{Истинная зависимость (Flow dependence / RAW --- Read-After-Write):} возникает, когда инструкция $S_1$ записывает данные в ячейку памяти, а инструкция $S_2$, выполняющаяся позже в исходном порядке программы, считывает данные из этой же ячейки. Истинная зависимость отражает передачу данных алгоритма и не может быть устранена переименованием переменных.
    \item \textbf{Антизависимость (Anti-dependence / WAR --- Write-After-Read):} возникает, когда инструкция $S_1$ считывает данные, а последующая инструкция $S_2$ перезаписывает их. Эта зависимость возникает из-за переиспользования памяти и часто может быть устранена путем приватизации переменных или расширения массивов.
    \item \textbf{Выходная зависимость (Output dependence / WAW --- Write-After-Write):} возникает, когда обе инструкции $S_1$ и $S_2$ записывают данные в одну и ту же ячейку памяти. Порядок этих записей критически важен, так как он определяет финальное значение ячейки.
\end{enumerate}

Для вложенных циклов особую роль играют циклически порожденные (переносимые) зависимости (loop-carried dependencies). Это зависимости, возникающие между операциями, выполняемыми на разных итерациях одного или нескольких объемлющих циклов. Именно они являются главным препятствием для свободного изменения порядка итераций и векторизации \cite{aho2007dragon}.

\subsection{Математическое моделирование: пространство итераций}

Для применения строгих математических методов произвольное гнездо из $n$ циклов рассматривается как целочисленное $n$-мерное пространство итераций. Каждой копии тела цикла ставится в соответствие вектор счётчиков $I = (i_1, i_2, \dots, i_n)^T$. Границы циклов задают систему линейных неравенств, определяющую множество допустимых точек --- выпуклый многогранник (полиэдр) в $\mathbb{Z}^n$ \cite{bondhugula2008pluto}.

В программах с регулярной структурой (где индексы массивов являются аффинными функциями от счетчиков циклов) компилятор строит решётчатое представление вычислений. Для описания направленности и "дальности" информационных зависимостей в этом пространстве вводятся векторы расстояний (distance vectors). Если истинная зависимость возникает между итерацией $I_{src}$ (генератор) и итерацией $I_{dst}$ (использование), то вектор расстояний определяется как разность: $D = I_{dst} - I_{src}$. 

Лексикографическая положительность векторов расстояний ($D \succ 0$) является фундаментальным условием корректности исходной последовательной программы: зависимость всегда направлена от более ранней итерации к более поздней. Цель компиляторных преобразований --- модифицировать пространство итераций (путем умножения на матрицы трансформации) таким образом, чтобы векторы расстояний приобрели свойства, необходимые для безопасного распараллеливания или тайлинга, не нарушая при этом условие $D' \succ 0$.

\subsection{Фундаментальные преобразования вложенных циклов}

В современном оптимизирующем трансляторе реализован ряд базовых полиэдральных и синтаксических преобразований, позволяющих адаптировать код под особенности целевой микроархитектуры \cite{allenKennedy2002}.

\subsubsection{Перестановка циклов (Loop Interchange)}
Данное преобразование изменяет порядок вложенности двух или более смежных циклов. Основная цель перестановки --- улучшение пространственной локальности. Стандарт Фортрана предписывает размещение элементов многомерных массивов в памяти по столбцам (column-major order). Это означает, что элементы $A(1, J)$ и $A(2, J)$ расположены в памяти физически смежно. Если в исходном коде внутренний цикл обходит второй индекс ($J$), программа будет обращаться к памяти с большим шагом, что приведет к катастрофическому количеству промахов кэш-памяти. Перестановка циклов так, чтобы внутренним стал цикл по первому индексу ($I$), восстанавливает последовательный доступ к памяти. Условием легальности перестановки является то, что после обмена соответствующих компонент ни один из векторов зависимостей не должен стать лексикографически отрицательным.

\subsubsection{Блочное разбиение (Loop Tiling / Loop Blocking)}
Тайлинг является одним из наиболее мощных инструментов оптимизации временной локальности. Пространство итераций многомерного цикла разбивается гиперплоскостями на компактные блоки (тайлы) заданного размера. Затем обход организуется таким образом, чтобы сначала выполнялись все вычисления внутри одного тайла, и лишь затем осуществлялся переход к следующему \cite{wolfLam1991}. 
Синтаксически это выражается в увеличении глубины гнезда циклов в два раза: внешние циклы (strip-mining) отвечают за перебор тайлов, а внутренние --- за выполнение итераций внутри выбранного тайла. Поскольку размер тайла подбирается таким образом, чтобы обрабатываемые им данные целиком помещались в кэш-память L1 или L2, тайлинг позволяет многократно переиспользовать загруженные данные, избегая их дорогостоящей выгрузки и повторной загрузки из оперативной памяти. Тайлинг допустим только в том случае, если все скалярные компоненты векторов зависимостей строго неотрицательны.

\subsubsection{Скашивание пространства итераций (Loop Skewing)}
Скашивание применяется в ситуациях, когда векторы зависимостей содержат отрицательные компоненты, препятствующие применению тайлинга или распараллеливания. Данное преобразование геометрически изменяет форму пространства итераций из ортогональной (прямоугольной) в параллелограммную \cite{lamport1974}.
Алгебраически скашивание эквивалентно умножению вектора итераций на нижнетреугольную унимодулярную матрицу. Скашивание с правильно подобранными коэффициентами сдвигает итерации внутреннего цикла пропорционально значениям внешнего цикла, что позволяет сделать отрицательные компоненты векторов зависимостей строго положительными. Таким образом, скашивание чаще всего выступает как подготовительный этап (enabling transformation) для последующей оптимизации.

\subsubsection{Метод гиперплоскостей (Loop Wavefront)}
В алгоритмах с множественными перекрестными зависимостями (например, при решении уравнений в частных производных неявными методами) циклы не могут быть распараллелены тривиальным образом. Метод гиперплоскостей, или метод волнового фронта, позволяет извлечь параллелизм из таких структур. Пространство итераций виртуально покрывается семейством параллельных гиперплоскостей. Направление продвижения по этим гиперплоскостям задает "волновой фронт" вычислений. Ключевое свойство заключается в том, что все итерации, лежащие строго в пределах одной гиперплоскости, гарантированно независимы друг от друга по данным и могут выполняться параллельно.

\subsubsection{Размотка циклов (Loop Unrolling)}
В отличие от вышеописанных трансформаций пространства итераций, размотка циклов является скалярной оптимизацией тела цикла. Она заключается в дублировании тела цикла несколько раз (с соответствующей корректировкой индексов) и уменьшении количества итераций во столько же раз. Размотка снижает накладные расходы на управление циклом (сравнение счетчика, условные переходы) и, что более важно, открывает широкие возможности для планировщика инструкций процессора (instruction scheduling) и автоматической векторизации, так как в развернутом теле цикла становится доступно больше независимых инструкций для загрузки исполнительных конвейеров суперскалярного процессора.

\subsubsection{Слияние и расщепление циклов (Loop Fusion and Fission)}
Слияние (fusion) объединяет несколько идущих подряд циклов с одинаковым пространством итераций в один общий цикл. Это позволяет сократить накладные расходы на инициализацию циклов и улучшает временную локальность, если оба цикла обращались к одним и тем же массивам данных. 
Расщепление (fission или distribution), напротив, разбивает один большой цикл на несколько независимых. Эта трансформация полезна для снижения давления на регистровый файл (register pressure), предотвращения вытеснения данных из кэша инструкций (instruction cache trashing) и изоляции фрагментов кода, содержащих циклически порожденные зависимости, от фрагментов, которые могут быть легко векторизованы.

\section{Методы ускорения гнёзд циклов итерационного типа}

\subsection{Определение и специфика гнёзд циклов итерационного типа}

Рассматриваемый в данной работе алгоритм фокусируется на специфическом, но чрезвычайно востребованном в научных расчетах классе структур --- гнёздах циклов итерационного типа. Подобная структура является канонической для большинства разностных схем, решателей систем линейных алгебраических уравнений итерационными методами и алгоритмов математического моделирования физических процессов.

Главной отличительной особенностью таких гнёзд является четкое функциональное разделение уровней вложенности. Внешний цикл (счетчик $t$ или $i_0$) отвечает за обход итераций самого численного алгоритма во времени (или шагов сходимости), тогда как группа внутренних циклов ($i_1, \dots, i_n$) реализует обход многомерной пространственной сетки по заданному фиксированному расчетному шаблону (stencil) \cite{metelica2024}.

Рассматриваются строго тесные гнёзда циклов, в которых генераторы (левые части операторов присваивания) имеют канонический вид $a_{i_1-c_1, \dots, i_n-c_n}$, а использования (операнды в правой части) --- $a_{i_1-d_1, \dots, i_n-d_n}$, где $c_k$ и $d_k$ --- небольшие целочисленные константы, определяющие "форму" расчетного шаблона. Важнейшим свойством является то, что счетчик внешнего итерационного цикла $i_0$ никогда явно не входит в индексные выражения массивов: предполагается, что алгоритм работает с несколькими буферами (или модифицирует матрицу "на месте" in-place), переходя от одного временного слоя к другому.

\subsection{Моделирование информационных зависимостей и вычисление векторов расстояний}

Первым шагом оптимизирующего компилятора является точное математическое моделирование информационных зависимостей. Для гнёзд итерационного типа размерности $n+1$ дуга информационной зависимости $(u, v)$ возникает между генератором и использованием, ссылающимися на один и тот же участок памяти.

Вектор расстояний для такой зависимости в пространстве $\mathbb{Z}^{n+1}$ формально определяется разностью векторов счетчиков циклов, при которых происходят зависимые обращения. Поскольку внешняя итерация всегда продвигается строго вперед во времени, первая координата вектора расстояний, соответствующая циклу $i_0$, всегда положительна (как правило, равна 1). Остальные координаты определяются разностями констант-смещений из индексных выражений. Таким образом, вектор расстояний для дуги $(u, v)$ строго задается формулой:
\begin{equation}
\label{eq:dist}
\operatorname{dist}(u, v) = (1, d_1 - c_1, d_2 - c_2, \dots, d_n - c_n).
\end{equation}

Наличие хотя бы одной отрицательной координаты (среди $d_k - c_k$) в любом из векторов расстояний сигнализирует компилятору о невозможности прямого применения блочного разбиения (тайлинга). Отрицательная компонента означает, что при наивном разбиении на блоки возникнет ситуация, когда блок, вычисление которого должно предшествовать другому блоку, будет лексикографически располагаться после него в пространстве итераций, что приведет к катастрофическому нарушению семантики программы.

\subsection{Алгоритм построения матрицы скашивания}

Для устранения отрицательных координат компилятор применяет алгоритм автоматического построения матрицы скашивания. Суть метода заключается в целенаправленном смещении (скосе) внутреннего пространства итераций относительно внешних циклов.

Для каждой дуги зависимости с отрицательной компонентой $\operatorname{dist}_m < 0$ (где $m \in \{1, \dots, n\}$) выполняется элементарное скашивание циклов-носителей этой зависимости относительно цикла $m$. Параметр скашивания $f$ выбирается равным модулю отрицательной координаты: $f = |\operatorname{dist}_m|$. Каждое такое элементарное скашивание задается нижнетреугольной матрицей $s$, где на главной диагонали расположены единицы, а элемент $s_{k, m} = f$ (для носителя $k$).

Если дуга имеет несколько отрицательных координат, формируется композиция матриц. Для всего множества дуг графа зависимостей компилятор собирает множество матриц единичных скашиваний $S$. Результирующая глобальная матрица скашивания $\mathrm{skew}$ вычисляется путем покомпонентного выбора максимальных элементов из всех матриц множества $S$:
\begin{equation}
\label{eq:skew_matrix}
\mathrm{skew}_{i,j} = \max_{s \in S} (s_{i,j}).
\end{equation}
Умножение исходного вектора счетчиков на полученную матрицу $\mathrm{skew}$ генерирует модифицированное пространство итераций, в котором гарантированно отсутствуют отрицательные компоненты векторов расстояний, что является достаточным условием для легального применения многомерного тайлинга.

\subsection{Блочное разбиение и оптимизация порядка обхода внутри тайла}

Основываясь на результатах скашивания, непрерывное пространство итераций разбивается ортогональными гиперплоскостями на многомерные блоки (тайлы) с заданными геометрическими размерами $d_1, d_2, \dots, d_n$. Данный этап принципиально важен для повышения временной локальности.

Однако классический тайлинг может быть существенно улучшен. Компилятор может модифицировать порядок обхода точек внутри уже сформированного тайла путем перестановки самого глубоко вложенного цикла, сделав его внешним (среди циклов, обходящих сам тайл). Геометрически это означает, что обход начинает производиться не по горизонтальным сечениям тайла, а по сечениям, параллельным его боковой грани \cite{metelica2024}.

Анализ баланса чтений памяти демонстрирует феноменальную эффективность данного решения. Обозначим количество чтений из оперативной памяти как $M$, из кэш-памяти L1 как $C$, а из сверхбыстрых регистров как $R$. Для разностного шаблона Гаусса-Зейделя стандартный обход по горизонтальным сечениям требует для большинства внутренних точек тайла одного чтения из регистров и трех чтений из кэша ($1R + 3C$). 

После предложенной компилятором перестановки обхода по боковым сечениям формула чтений трансформируется в $2R + 2C$. Удвоение количества чтений непосредственно из аппаратных регистров процессора, доступ к которым осуществляется без задержек тактов, приводит к существенному снижению нагрузки на подсистему кэш-памяти. В масштабах миллиардов итераций это структурное изменение обеспечивает значительное дополнительное ускорение, превосходящее результаты классических систем полиэдральной оптимизации (таких как PLUTO) до 1.4 раза.

\subsection{Метод гиперплоскостей и распараллеливание тайлов}

Финальным этапом глобальной структурной оптимизации является автоматическая генерация многопоточного кода для многоядерных архитектур. После разбиения пространства на блоки формируется фактор-граф по тайлам, вершинами которого выступают целые вычислительные блоки. 

Компилятор применяет метод гиперплоскостей (wavefront) к этому фактор-графу, покрывая его семейством параллельных многомерных гиперплоскостей с единичным вектором нормали $N = (1, 1, \dots, 1)$. Выбор единичной нормали максимизирует количество тайлов, попадающих на одну плоскость, обеспечивая хорошую балансировку нагрузки (load balancing) между вычислительными потоками.

Фундаментальная теорема эквивалентности гласит, что после корректного применения скашивания (обеспечивающего неотрицательность векторов расстояний) и тайлинга, любые два тайла, лежащие строго на одной гиперплоскости с нормалью $(1, 1, \dots, 1)$, попарно не связаны никакими информационными зависимостями. 

Следовательно, компилятор имеет полное математическое право сгенерировать параллельный код (например, используя директивы `!$OMP PARALLEL DO` в транслированном Фортран-коде) для цикла, обходящего тайлы в пределах одной плоскости фронта. Это преобразование является строго эквивалентным и гарантирует отсутствие состояний гонки за данными (data races).

\subsection{Аналитическая модель выбора оптимальных размеров тайлов}

Проблема автоматического определения параметров блочного разбиения (tile size selection) решается в компиляторе посредством встроенной аналитической модели, связывающей геометрию тайла с физическими характеристиками целевого процессора \cite{metelica2024}. Модель оценивает объем памяти $V$, необходимый для хранения множества элементов массива, которые интенсивно переиспользуются в процессе вычисления одного шага внутри тайла.

Для двумерного расчетного шаблона минимизация промахов в L1 кэш (L1 cache misses) требует, чтобы совокупный рабочий набор данных трех смежных сечений тайла непрерывно и целиком помещался в кэш-памяти. Обозначим физическую емкость кэша L1, доступную одному ядру, как $|L1|$. Количество уникальных элементов массива, вовлеченных в вычисление одного внутреннего сечения с учетом краевых точек шаблона, оценивается как:
\begin{equation}
\label{eq:volume}
V \le d_1 d_2 + 2(d_1 + d_2) = (d_1 + 2)(d_2 + 2) - 4.
\end{equation}
Отсюда вытекает фундаментальное физическое ограничение для оптимизирующего компилятора:
\begin{equation}
\label{eq:l1_limit}
(d_1 + 2)(d_2 + 2) - 4 \le |L1|.
\end{equation}
Математически известно, что при фиксированной площади сечения ($d_1 \times d_2$) его периметр, отвечающий за подкачку граничных значений из более медленных уровней памяти, достигает минимума в форме квадрата, то есть при $d_1 = d_2$. Разрешая неравенство относительно симметричного тайла, компилятор получает аналитическое выражение для теоретически оптимального размера:
\begin{equation}
\label{eq:opt_tile}
d_1 = d_2 \approx \sqrt{|L1| + 4} - 2.
\end{equation}
Для иллюстрации, при объеме L1 кэша в 32 Кбайт, способном вместить 4096 8-байтовых чисел с плавающей точкой двойной точности (тип REAL8), оптимальный размер тайла составит $d_1 = d_2 \approx \sqrt{4100} - 2 \approx 62$. Эмпирические испытания подтверждают, что максимальное ускорение на многоядерных архитектурах достигается именно в окрестности предсказанных размеров, что делает данную аналитическую модель надежным инструментом для автоматической настройки параметров (auto-tuning) компилятора.

\subsection{Дополнительные скалярные оптимизации линейных участков}

Для раскрытия полного потенциала оптимизированного кода, агрессивные глобальные цикловые трансформации пространства итераций должны сопровождаться эффективными локальными скалярными оптимизациями в рамках сгенерированных блоков. 

Изменение топологии гнезда циклов неминуемо усложняет математическую структуру индексных выражений. В тела циклов внедряются конструкции, включающие умножения счетчиков тайлов на константы $d_i$, вычисления минимальных/максимальных границ и сложные алгебраические смещения. Для минимизации накладных расходов от этих вычислений компилятор применяет:
\begin{itemize}
    \item \textbf{Вынесение инвариантных выражений (Loop Invariant Code Motion):} Идентификация и принудительное вынесение сложных подвыражений, операнды которых константны в рамках конкретного цикла, за пределы его заголовка. Это предотвращает их многократное перевычисление на каждой итерации.
    \item \textbf{Линеаризацию и алгебраическое упрощение выражений (Constant Folding \& Strength Reduction):} Приведение сложных многомерных аффинных индексных выражений к каноническому, максимально упрощенному виду. Производится замена медленных целочисленных умножений на индуктивные сложения и сдвиги.
\end{itemize}
Комплексное сочетание глобальной реорганизации пространства итераций с локальной микрооптимизацией скалярных операций генерирует исключительно высокопроизводительный машинный код, обеспечивая дополнительное ускорение (помимо распараллеливания) в 1.5--1.8 раза.

\anonsection{ЗАКЛЮЧЕНИЕ}

В рамках настоящей научно-исследовательской работы было проведено глубокое теоретическое обоснование и проработка методов автоматической оптимизации многомерных гнёзд циклов. Целью исследования являлось формирование надежной алгоритмической и математической базы для создания архитектуры современного оптимизирующего компилятора подмножества языка Фортран, ориентированного на достижение пиковой производительности в научных численных расчетах.

В ходе выполнения исследования получены следующие фундаментальные результаты:
\begin{enumerate}
    \item Проведен комплексный анализ теоретических основ цикловых оптимизаций. Доказана принципиальная роль статического анализа потока данных, выявления информационных зависимостей и формального геометрического моделирования программ с использованием аппарата пространств итераций.
    \item Детально исследована и математически обоснована концепция локальности данных (пространственной и временной) применительно к современным процессорным архитектурам с иерархической системой кэш-памяти.
    \item Разработан, теоретически обоснован и адаптирован комплексный многостадийный алгоритм ускорения гнёзд циклов итерационного типа. Установлено, что вычисление векторов расстояний информационных зависимостей позволяет оптимизатору автоматически синтезировать матрицу скашивания (loop skewing matrix), устраняющую отрицательные компоненты и математически разрешающую последующее применение блочного разбиения (тайлинга).
    \item Доказана критическая важность микрооптимизации обхода внутри самого тайла. Математический анализ потока чтений показал, что целенаправленная перестановка циклов, обеспечивающая обход по боковым сечениям тайла, способна радикально увеличить долю обращений процессора напрямую к сверхбыстрым аппаратным регистрам. Эта трансформация снижает нагрузку на кэш-память L1 и позволяет превзойти по производительности аналогичные решения в существующих полиэдральных оптимизаторах.
    \item Сформулировано и приведено строгое доказательство теоремы об информационной независимости тайлов, лежащих в пределах одной геометрической гиперплоскости с единичным вектором нормали. Это доказательство служит надежным теоретическим фундаментом для автоматического и абсолютно безопасного (без риска гонок за данными) крупноблочного распараллеливания вычислений на системы с общей памятью с использованием стандарта OpenMP.
    \item Разработана элегантная аналитическая модель, связывающая теоретически оптимальные геометрические размеры тайла с физическими аппаратными ограничениями микропроцессора. Математически выведено, что для минимизации промахов кэш-памяти при вычислении регулярных пространственных шаблонов, оптимальный размер квадратного сечения тайла находится в прямой зависимости от емкости L1-кэша. Это открывает возможность для создания механизмов адаптивного автоматического тюнинга (auto-tuning) в компиляторе.
\end{enumerate}

Подводя итог, можно заключить, что разработанная теоретическая база представляет собой концептуально завершенную, строгую математическую систему трансформаций. Применение описанной последовательности --- от анализа зависимостей и скашивания до тайлинга, перестановки, линеаризации выражений и параллелизации методом гиперплоскостей --- формирует полноценный конвейер (pipeline) для интеллектуального компилятора подмножества Фортрана нового поколения. Реализация предложенного теоретического базиса в программном коде транслятора гарантирует автоматическое многократное ускорение критически важных вычислительных модулей, что составляет весомый научно-практический вклад в инструментарий высокопроизводительных вычислений.
\renewcommand\refname{СПИСОК ИСПОЛЬЗОВАННЫХ ИСТОЧНИКОВ}
\clearpage
{\catcode`"\active\def"{\relax}
\addcontentsline{toc}{section}{\protect\numberline{}\refname}%
\printbibliography
}

\end{document}
